% Part: cheat-sheets
% Chapter: first-order-logic

\documentclass[../../../../include/open-logic-chapter]{subfiles}

\begin{document}

\olchapter{chs}{fol}{First-Order Logic}

\section{Semantics}

\begin{defn}[!!^{structure}s]
\Article{structure} \emph{!!{structure}}~$\Struct M$, for a language
$\Lang{L}$ of first-order logic consists of the following elements:
\begin{enumerate}
\item \emph{Domain:} a non-empty set, $\Domain M$
\item \emph{Interpretation of !!{constant}s:} for each !!{constant}~$c$ of
  $\Lang{L}$, !!a{element} $\Assign{c}{M} \in \Domain M$
\item \emph{Interpretation of !!{predicate}s:} for each $n$-place
  !!{predicate}~$R$ of $\Lang{L}$ (other than $\eq$), an $n$-place
  relation $\Assign{R}{M} \subseteq \Domain{M}^n$
\tagitem{fnTerms}{\emph{Interpretation of !!{function}s:} for each $n$-place
  !!{function}~$f$ of $\Lang{L}$, an $n$-place function $\Assign{f}{M}
  \colon \Domain{M}^n \to \Domain{M}$}{}
\end{enumerate}
\end{defn}

\begin{defn}[!!^{value} of Terms]
If $t$ is a term of the language~$\Lang L$, $\Struct M$ is a
!!{structure} for~$\Lang L$, and $s$ is a !!{variable} assignment
for~$\Struct M$, the \emph{!!{value}}~$\Value{t}{M}[s]$ is defined as
follows:
\begin{enumerate}
\item \indcase{t}{c}{$\Value{\indfrm}{M}[s] = \Assign{\indcomplex}{M}$.}
\item \indcase{t}{x}{$\Value{\indfrm}{M}[s] = s(\indcomplex)$.}
\tagitem{fnTerms}{\indcase{t}{\Atom{f}{t_1, \ldots, t_n}}{
\[
\Value{\indfrm}{M}[s] = \Assign{f}{M}(\Value{t_1}{M}[s], \ldots,
\Value{t_n}{M}[s]).
\]}
}{}
\end{enumerate}
\end{defn}

\begin{defn}[$x$-Variant]
If $s$ is a !!{variable} assignment for a !!{structure}~$\Struct M$, then any
!!{variable} assignment $s'$ for $\Struct M$ which differs from $s$ at most
in what it assigns to $x$ is called an \emph{$x$-variant} of~$s$.  If
$s'$ is an $x$-variant of $s$ we write $\varAssign{s'}{s}{x}$.
\end{defn}

\begin{defn}[Satisfaction]
\ollabel{defn:satisfaction}
Satisfaction of a !!{formula}~$!A$ in a !!{structure}~$\Struct M$
relative to a !!{variable} assignment~$s$, in symbols:
$\Sat{M}{!A}[s]$, is defined recursively as follows. (We write
$\Sat/{M}{!A}[s]$ to mean ``not $\Sat{M}{!A}[s]$.'')
\begin{enumerate}
\tagitem{prvFalse}{%
  \indcase{!A}{\lfalse}{$\Sat/{M}{\indfrm}[s]$.}}{}

\tagitem{prvTrue}{%
  \indcase{!A}{\ltrue}{$\Sat{M}{\indfrm}[s]$.}}{}

\item \indcase{!A}{\Atom{R}{t_1, \dots, t_n}}{$\Sat{M}{\indfrm}[s]$
  iff $\langle \Value{t_1}{M}[s], \dots, \Value{t_n}{M}[s] \rangle \in
  \Assign{R}{M}$.}

\item \indcase{!A}{\eq[t_1][t_2]}{$\Sat{M}{\indfrm}[s]$ iff
  $\Value{t_1}{M}[s] = \Value{t_2}{M}[s]$.}

\tagitem{prvNot}{%
  \indcase{!A}{\lnot !B}{$\Sat{M}{\indfrm}[s]$ iff
    $\Sat/{M}{!B}[s]$.}}{}

\tagitem{prvAnd}{%
  \indcase{!A}{(!B \land !C)}{$\Sat{M}{\indfrm}[s]$ iff $\Sat{M}{!B}[s]$
    and $\Sat{M}{!C}[s]$.}}{}

\tagitem{prvOr}{%
  \indcase{!A}{(!B \lor !C)}{$\Sat{M}{\indfrm}[s]$ iff
    $\Sat{M}{!A}[s]$ or $\Sat{M}{!B}[s]$ (or both).}}{}

\tagitem{prvIf}{%
  \indcase{!A}{(!B \lif !C)}{$\Sat{M}{\indfrm}[s]$ iff $\Sat/{M}{!B}[s]$
    or $\Sat{M}{!C}[s]$ (or both).}}{}

\tagitem{prvIff}{%
  \indcase{!A}{(!B \liff !C)}{$\Sat{M}{\indfrm}[s]$ iff either both
    $\Sat{M}{!B}[s]$ and $\Sat{M}{!C}[s]$, or neither $\Sat{M}{!B}[s]$
    nor $\Sat{M}{!C}[s]$.}}{}

\tagitem{prvAll}{%
  \indcase{!A}{\lforall[x][!B]}{$\Sat{M}{\indfrm}[s]$ iff for every
    $x$-variant~$s'$ of $s$, $\Sat{M}{!B}[s']$.}}{}

\tagitem{prvEx}{%
  \indcase{!A}{\lexists[x][!B]}{$\Sat{M}{\indfrm}[s]$ iff there is an
    $x$-variant~$s'$ of $s$ so that $\Sat{M}{!B}[s']$.}}{}
\end{enumerate}
\end{defn}

\begin{defn}[Validity]
A sentence $!A$ is \emph{valid}, $\Entails !A$, iff $\Sat{M}{!A}$ for every
!!{structure}~$\Struct M$.
\end{defn}

\begin{defn}[Entailment]
A set of sentences~$\Gamma$ \emph{entails} a sentence~$!A$, $\Gamma
\Entails !A$, iff for every !!{structure}~$\Struct M$ with
$\Sat{M}{\Gamma}$, $\Sat{M}{!A}$.
\end{defn}


\begin{defn}[Satisfiability]
A set of sentences~$\Gamma$ is \emph{satisfiable} if $\Sat{M}{\Gamma}$
for some !!{structure}~$\Struct M$.  If $\Gamma$ is not satisfiable it is
called \emph{unsatisfiable}.
\end{defn}

\section{Tableaux}

\subsection{Rules for $\lforall$}

\begin{defish}
\AxiomC{\sFmla{\True}{\lforall[x][!A(x)]}}
\RightLabel{\TRule{\True}{\forall}}
\UnaryInfC{\sFmla{\True}{!A(t)}}
\DisplayProof
\hfill
\AxiomC{\sFmla{\False}{\lforall[x][!A(x)]}}
\RightLabel{\TRule{\False}{\lforall}}
\UnaryInfC{\sFmla{\False}{!A(a)}}
\DisplayProof
\end{defish}

In \TRule{\True}{\lforall}, $t$ is a closed term (i.e., one without
variables). In \TRule{\False}{\lforall}, $a$~is !!a{constant} which must
not occur anywhere in the branch above \TRule{\False}{\lforall}
rule. We call $a$ the \emph{eigenvariable} of the \TRule{\False}{\forall}
inference.

\subsection{Rules for $\lexists$}

\begin{defish}
\AxiomC{\sFmla{\True}{\lexists[x][!A(x)]}}
\RightLabel{\TRule{\True}{\lexists}}
\UnaryInfC{\sFmla{\True}{!A(a)}}
\DisplayProof
\hfill
\AxiomC{\sFmla{\False}{\lexists[x][!A(x)]}}
\RightLabel{\TRule{\False}{\lexists}}
\UnaryInfC{\sFmla{\False}{!A(t)}}
\DisplayProof
\end{defish}

\begin{defish}
\AxiomC{}
\RightLabel{$\eq$}
\UnaryInfC{\sFmla{\True}{\eq[t][t]}}
\DisplayProof
\hfill
\AxiomC{\sFmla{\True}{\eq[t_1][t_2]}}
\noLine
\UnaryInfC{\sFmla{\True}{!A(t_1)}}
\RightLabel{$\TRule{\True}{\eq}$}
\UnaryInfC{\sFmla{\True}{!A(t_2)}}
\DisplayProof
\hfill
\AxiomC{\sFmla{\True}{\eq[t_1][t_2]}}
\noLine
\UnaryInfC{\sFmla{\False}{!A(t_1)}}
\RightLabel{$\TRule{\True}{\eq}$}
\UnaryInfC{\sFmla{\False}{!A(t_2)}}
\DisplayProof
\end{defish}

\OLEndChapterHook

\end{document}
